% Options for packages loaded elsewhere
% Options for packages loaded elsewhere
\PassOptionsToPackage{unicode}{hyperref}
\PassOptionsToPackage{hyphens}{url}
\PassOptionsToPackage{dvipsnames,svgnames,x11names}{xcolor}
%
\documentclass[
  12pt,
  letterpaper,
  DIV=11,
  numbers=noendperiod]{scrartcl}
\usepackage{xcolor}
\usepackage{amsmath,amssymb}
\setcounter{secnumdepth}{-\maxdimen} % remove section numbering
\usepackage{iftex}
\ifPDFTeX
  \usepackage[T1]{fontenc}
  \usepackage[utf8]{inputenc}
  \usepackage{textcomp} % provide euro and other symbols
\else % if luatex or xetex
  \usepackage{unicode-math} % this also loads fontspec
  \defaultfontfeatures{Scale=MatchLowercase}
  \defaultfontfeatures[\rmfamily]{Ligatures=TeX,Scale=1}
\fi
\usepackage{lmodern}
\ifPDFTeX\else
  % xetex/luatex font selection
  \setmainfont[]{Times New Roman}
\fi
% Use upquote if available, for straight quotes in verbatim environments
\IfFileExists{upquote.sty}{\usepackage{upquote}}{}
\IfFileExists{microtype.sty}{% use microtype if available
  \usepackage[]{microtype}
  \UseMicrotypeSet[protrusion]{basicmath} % disable protrusion for tt fonts
}{}
\makeatletter
\@ifundefined{KOMAClassName}{% if non-KOMA class
  \IfFileExists{parskip.sty}{%
    \usepackage{parskip}
  }{% else
    \setlength{\parindent}{0pt}
    \setlength{\parskip}{6pt plus 2pt minus 1pt}}
}{% if KOMA class
  \KOMAoptions{parskip=half}}
\makeatother
% Make \paragraph and \subparagraph free-standing
\makeatletter
\ifx\paragraph\undefined\else
  \let\oldparagraph\paragraph
  \renewcommand{\paragraph}{
    \@ifstar
      \xxxParagraphStar
      \xxxParagraphNoStar
  }
  \newcommand{\xxxParagraphStar}[1]{\oldparagraph*{#1}\mbox{}}
  \newcommand{\xxxParagraphNoStar}[1]{\oldparagraph{#1}\mbox{}}
\fi
\ifx\subparagraph\undefined\else
  \let\oldsubparagraph\subparagraph
  \renewcommand{\subparagraph}{
    \@ifstar
      \xxxSubParagraphStar
      \xxxSubParagraphNoStar
  }
  \newcommand{\xxxSubParagraphStar}[1]{\oldsubparagraph*{#1}\mbox{}}
  \newcommand{\xxxSubParagraphNoStar}[1]{\oldsubparagraph{#1}\mbox{}}
\fi
\makeatother


\usepackage{longtable,booktabs,array}
\newcounter{none} % for unnumbered tables
\usepackage{calc} % for calculating minipage widths
% Correct order of tables after \paragraph or \subparagraph
\usepackage{etoolbox}
\makeatletter
\patchcmd\longtable{\par}{\if@noskipsec\mbox{}\fi\par}{}{}
\makeatother
% Allow footnotes in longtable head/foot
\IfFileExists{footnotehyper.sty}{\usepackage{footnotehyper}}{\usepackage{footnote}}
\makesavenoteenv{longtable}
\usepackage{graphicx}
\makeatletter
\newsavebox\pandoc@box
\newcommand*\pandocbounded[1]{% scales image to fit in text height/width
  \sbox\pandoc@box{#1}%
  \Gscale@div\@tempa{\textheight}{\dimexpr\ht\pandoc@box+\dp\pandoc@box\relax}%
  \Gscale@div\@tempb{\linewidth}{\wd\pandoc@box}%
  \ifdim\@tempb\p@<\@tempa\p@\let\@tempa\@tempb\fi% select the smaller of both
  \ifdim\@tempa\p@<\p@\scalebox{\@tempa}{\usebox\pandoc@box}%
  \else\usebox{\pandoc@box}%
  \fi%
}
% Set default figure placement to htbp
\def\fps@figure{htbp}
\makeatother

\ifLuaTeX
  \usepackage{luacolor}
  \usepackage[soul]{lua-ul}
\else
  \usepackage{soul}
\fi

% definitions for citeproc citations
\NewDocumentCommand\citeproctext{}{}
\NewDocumentCommand\citeproc{mm}{%
  \begingroup\def\citeproctext{#2}\cite{#1}\endgroup}
\makeatletter
 % allow citations to break across lines
 \let\@cite@ofmt\@firstofone
 % avoid brackets around text for \cite:
 \def\@biblabel#1{}
 \def\@cite#1#2{{#1\if@tempswa , #2\fi}}
\makeatother
\newlength{\cslhangindent}
\setlength{\cslhangindent}{1.5em}
\newlength{\csllabelwidth}
\setlength{\csllabelwidth}{3em}
\newenvironment{CSLReferences}[2] % #1 hanging-indent, #2 entry-spacing
 {\begin{list}{}{%
  \setlength{\itemindent}{0pt}
  \setlength{\leftmargin}{0pt}
  \setlength{\parsep}{0pt}
  % turn on hanging indent if param 1 is 1
  \ifodd #1
   \setlength{\leftmargin}{\cslhangindent}
   \setlength{\itemindent}{-1\cslhangindent}
  \fi
  % set entry spacing
  \setlength{\itemsep}{#2\baselineskip}}}
 {\end{list}}
\usepackage{calc}
\newcommand{\CSLBlock}[1]{\hfill\break\parbox[t]{\linewidth}{\strut\ignorespaces#1\strut}}
\newcommand{\CSLLeftMargin}[1]{\parbox[t]{\csllabelwidth}{\strut#1\strut}}
\newcommand{\CSLRightInline}[1]{\parbox[t]{\linewidth - \csllabelwidth}{\strut#1\strut}}
\newcommand{\CSLIndent}[1]{\hspace{\cslhangindent}#1}



\setlength{\emergencystretch}{3em} % prevent overfull lines

\providecommand{\tightlist}{%
  \setlength{\itemsep}{0pt}\setlength{\parskip}{0pt}}



 


\KOMAoption{captions}{tableheading}
\makeatletter
\@ifpackageloaded{caption}{}{\usepackage{caption}}
\AtBeginDocument{%
\ifdefined\contentsname
  \renewcommand*\contentsname{Table of contents}
\else
  \newcommand\contentsname{Table of contents}
\fi
\ifdefined\listfigurename
  \renewcommand*\listfigurename{List of Figures}
\else
  \newcommand\listfigurename{List of Figures}
\fi
\ifdefined\listtablename
  \renewcommand*\listtablename{List of Tables}
\else
  \newcommand\listtablename{List of Tables}
\fi
\ifdefined\figurename
  \renewcommand*\figurename{Figure}
\else
  \newcommand\figurename{Figure}
\fi
\ifdefined\tablename
  \renewcommand*\tablename{Table}
\else
  \newcommand\tablename{Table}
\fi
}
\@ifpackageloaded{float}{}{\usepackage{float}}
\floatstyle{ruled}
\@ifundefined{c@chapter}{\newfloat{codelisting}{h}{lop}}{\newfloat{codelisting}{h}{lop}[chapter]}
\floatname{codelisting}{Listing}
\newcommand*\listoflistings{\listof{codelisting}{List of Listings}}
\makeatother
\makeatletter
\makeatother
\makeatletter
\@ifpackageloaded{caption}{}{\usepackage{caption}}
\@ifpackageloaded{subcaption}{}{\usepackage{subcaption}}
\makeatother
\usepackage{bookmark}
\IfFileExists{xurl.sty}{\usepackage{xurl}}{} % add URL line breaks if available
\urlstyle{same}
\hypersetup{
  pdftitle={Protocol},
  pdfauthor={Mikkel Errboe Steentoft},
  colorlinks=true,
  linkcolor={blue},
  filecolor={Maroon},
  citecolor={Blue},
  urlcolor={Blue},
  pdfcreator={LaTeX via pandoc}}


\title{Protocol}
\author{Mikkel Errboe Steentoft}
\date{}
\begin{document}
\maketitle


\section{Continous glucose monitoring for metabolic profiling in
prediabetes}\label{continous-glucose-monitoring-for-metabolic-profiling-in-prediabetes}

\subsection{Background}\label{background}

Prediabetes is a state of elevated blood glucose levels below the
threshold for type 2 diabetes and is associated with increased risk of
diabetes {[}1{]}, cardiovascular disease (CVD) and mortality {[}2{]}.
According to the international Diabetes Federation, 537 million adults
worldwide were living with prediabetes in 2021, which is projected to
increase to 783 million by 2045 {[}3{]}, highlighting a global need for
prevention strategies in prediabetes. Prediabetes is currently defined
by different glycemic thresholds; elevated glycated hemoglobin (HbA1c)
level of 39-47 mmol/l {[}4{]} or 42-47 mmol/l {[}5{]}, elevated fasting
glucose levels (IFG) or elevated 2-hour-glucose (2hPG) after an oral
glucose tolerance test (OGTT). While current diagnostic criteria for
prediabetes capture the degree of dysglycemia at a single time-point,
they offer limited insight into underlying pathophysiology. Individuals
meeting the same glycemic criteria for prediabetes may differ
substantially in their underlying metabolic drivers and risk
trajectories. This includes distinct characteristics of insulin
resistance, beta-cell function and body fat distribution along with
varying risks of developing diabetes and related complications {[}6{]}.
Given this substantial heterogeneity among individuals with prediabetes,
effective prevention requires the ability to distinguish high-risk
individuals who could benefit from early, intensive intervention from
low-risk individuals who require minimal clinical attention {[}7{]}.
However, as also recognized in these studies, assessment of the
underlying metabolic drivers currently relies on complex tools, such as
the OGTT for measuring post-load glucose regulation, or MRI and
ultrasonography for the assessment of body fat distribution {[}8,
9{]}.An important step towards precision prevention is therefore to
develop more clinically accessible tools capable of capturing underlying
metabolic dimensions.

Continuous glucose monitoring (CGM) offers a possibility to assess
detailed and dynamic indices of glucose regulation under free-living
conditions, and is already widely used in the management of type 1 and
type 2 diabetes {[}10{]}. Emerging evidence suggests that the use of CGM
could also be beneficial in the management of individuals with
prediabetes {[}11{]}. Recent studies including large-scale machine
learning models have shown that CGM- and wearable-derived measures can
predict future glycemic deterioration, incident type 2 diabetes,
cardiovascular outcomes {[}12{]} and insulin resistance {[}13{]}, often
outperforming HbA1c. While these approaches have the potential to
provide important information for risk stratification, they provide
limited insight into the distinction between insulin resistance and
beta-cell function - a distinction that may be important not only for
identifying who is at risk, but also for determining which type of
preventative strategy, such as intensive lifestyle interventions or
early pharmacological treatment, is most suitable for a given
individual{[}6{]}. In this context, it remains unclear whether CGM- and
wearable-derived information in free-living conditions can help identify
underlying metabolic dimensions, such as insulin resistance and
beta-cell function, among individuals with prediabetes, and how this
could inform individualized prevention strategies.

\subsection{Aims}\label{aims}

The overall aim of this PhD project is to assess whether continuous
glucose monitoring, combined with other free-living data, can help
characterize underlying metabolic dimensions, including insulin
resistance and beta-cell function, among individuals with prediabetes.
Through different data sources, the project spans from observational,
cross-sectional and longitudinal studies to detailed physiological
characterization.

The project will be comprised of three studies:

{\def\LTcaptype{none} % do not increment counter
\begin{longtable}[]{@{}
  >{\raggedright\arraybackslash}p{(\linewidth - 6\tabcolsep) * \real{0.3541}}
  >{\raggedright\arraybackslash}p{(\linewidth - 6\tabcolsep) * \real{0.4027}}
  >{\raggedright\arraybackslash}p{(\linewidth - 6\tabcolsep) * \real{0.1732}}
  >{\raggedright\arraybackslash}p{(\linewidth - 6\tabcolsep) * \real{0.0661}}@{}}
\toprule\noalign{}
\endhead
\bottomrule\noalign{}
\endlastfoot
Study/Aim & Data source & Population & Study design \\
\ul{\textbf{Study 1}}

Describe how CGM-derived glycemic patterns relate to free-living
lifestyle data (physical activity, sleep, diet) and standard clinical
measures among individuals with prediabetes. &
\begin{minipage}[t]{\linewidth}\raggedright
Baseline data from the DP-NEXY study:
\href{https://dp-next.github.io/wp3.html}{DP-NEXT, Work Package 3:
Heterogeneity}, a Danish study that will characterize heterogeneity
among individuals with prediabetes.

\hfill\break
\strut
\end{minipage} & Danish adults aged between 40-50 years with an HbA1c
measurement between 42-47 mmol/l & Cross-sectional \\
\ul{\textbf{Study 2}}

Investigate how CGM-derived patterns and related lifestyle
characteristics associates with long-term outcomes, including type 2
diabetes, cardiovascular disease and mortality. & Baseline and follow-up
data from \href{https://www.demaastrichtstudie.nl/research}{The
Maastricht Study}, a Dutch phenotyping study focusing on the etiology of
type 2 diabetes. & Dutch adults aged between 40-75 years enriched with
individuals with type 2 diabetes. & Longitudinal \\
\ul{\textbf{Study 3}}

Characterize underlying metabolic dimensions, including insulin
resistance, beta-cell function and substrate metabolism and how they
relate to CGM response patterns. & Exploratory physiological study using
whole-room Indirect calorimetry at steno diabetes center Aarhus & 5-10
participants from the \href{https://dp-next.github.io}{The DP-NEXT
study} & Exploratory observational study \\
\end{longtable}
}

\protect\phantomsection\label{refs}
\begin{CSLReferences}{0}{1}
\bibitem[\citeproctext]{ref-tabuxe1k2012}
\CSLLeftMargin{1. }%
\CSLRightInline{Tabák AG, Herder C, Rathmann W, Brunner EJ, Kivimäki M
(2012) Prediabetes: A high-risk state for developing diabetes. Lancet
379(9833):2279--2290.
\url{https://doi.org/10.1016/S0140-6736(12)60283-9}}

\bibitem[\citeproctext]{ref-vistisen2018}
\CSLLeftMargin{2. }%
\CSLRightInline{Vistisen D, Witte DR, Brunner EJ, et al (2018) Risk of
Cardiovascular Disease and Death in Individuals With Prediabetes Defined
by Different Criteria: The Whitehall II Study. Diabetes Care
41(4):899--906. \url{https://doi.org/10.2337/dc17-2530}}

\bibitem[\citeproctext]{ref-sun2022}
\CSLLeftMargin{3. }%
\CSLRightInline{Sun H, Saeedi P, Karuranga S, et al (2022) IDF Diabetes
Atlas: Global, regional and country-level diabetes prevalence estimates
for 2021 and projections for 2045. Diabetes Research and Clinical
Practice 183:109119.
\url{https://doi.org/10.1016/j.diabres.2021.109119}}

\bibitem[\citeproctext]{ref-americandiabetesassociationprofessionalpracticecommittee2025}
\CSLLeftMargin{4. }%
\CSLRightInline{American Diabetes Association Professional Practice
Committee (2025) 2. Diagnosis and Classification of Diabetes: Standards
of Care in Diabetes-2025. Diabetes Care 48(1 Suppl 1):S27--S49.
\url{https://doi.org/10.2337/dc25-S002}}

\bibitem[\citeproctext]{ref-internat2009}
\CSLLeftMargin{5. }%
\CSLRightInline{(2009) International expert committee report on the role
of the A1C assay in the diagnosis of diabetes. Diabetes Care
32(7):1327--1334. \url{https://doi.org/10.2337/dc09-9033}}

\bibitem[\citeproctext]{ref-wagner2026}
\CSLLeftMargin{6. }%
\CSLRightInline{Wagner R, Selvin E, Sehgal R, et al (2026) Beyond
Glucose-Rethinking Prediabetes for Precision Prevention. Diabetes Care
49(2):226--235. \url{https://doi.org/10.2337/dci25-0054}}

\bibitem[\citeproctext]{ref-blond2023}
\CSLLeftMargin{7. }%
\CSLRightInline{Blond MB, Færch K, Herder C, Ziegler D, Stehouwer CDA
(2023) The prediabetes conundrum: striking the balance between risk and
resources. Diabetologia 66(6):1016--1023.
\url{https://doi.org/10.1007/s00125-023-05890-y}}

\bibitem[\citeproctext]{ref-wagner2021}
\CSLLeftMargin{8. }%
\CSLRightInline{Wagner R, Heni M, Tabák AG, et al (2021)
Pathophysiology-based subphenotyping of individuals at elevated risk for
type 2 diabetes. Nature Medicine 27(1):49--57.
\url{https://doi.org/10.1038/s41591-020-1116-9}}

\bibitem[\citeproctext]{ref-stroebel2025}
\CSLLeftMargin{9. }%
\CSLRightInline{Stroebel BM, Gadgil M, Lewis KA, Longoria KD, Zhang L,
Flowers E (2025) Body composition differentiates prediabetes phenotype
clusters in the diabetes prevention program study. The Journal of
Clinical Endocrinology \& Metabolism 110(11):e3665--e3672.
\url{https://doi.org/10.1210/clinem/dgaf163}}

\bibitem[\citeproctext]{ref-battelino2025}
\CSLLeftMargin{10. }%
\CSLRightInline{Battelino T, Lalic N, Hussain S, et al (2025) The use of
continuous glucose monitoring in people living with obesity,
intermediate hyperglycemia or type 2 diabetes. Diabetes Research and
Clinical Practice 223:112111.
\url{https://doi.org/10.1016/j.diabres.2025.112111}}

\bibitem[\citeproctext]{ref-zahalka2024}
\CSLLeftMargin{11. }%
\CSLRightInline{Zahalka SJ, Galindo RJ, Shah VN, Low Wang CC (2024)
Continuous Glucose Monitoring for Prediabetes: What Are the Best
Metrics? Journal of Diabetes Science and Technology 18(4):835--846.
\url{https://doi.org/10.1177/19322968241242487}}

\bibitem[\citeproctext]{ref-lutsker2026}
\CSLLeftMargin{12. }%
\CSLRightInline{Lutsker G, Sapir G, Shilo S, et al (2026) A foundation
model for continuous glucose monitoring data. Nature 650(8103):978--986.
\url{https://doi.org/10.1038/s41586-025-09925-9}}

\bibitem[\citeproctext]{ref-insulin}
\CSLLeftMargin{13. }%
\CSLRightInline{\href{https://www.nature.com/articles/s41586-026-10179-2}{Insulin
resistance prediction from wearables and routine blood biomarkers
\textbar{} nature}}

\end{CSLReferences}




\end{document}
